\documentclass[10pt,a4paper]{amsart}
\usepackage[margin=19mm]{geometry}
\usepackage{amsmath,amssymb,mathtools}
\usepackage[scr=rsfs,cal=euler]{mathalfa}
\usepackage{xfrac}
\usepackage[unicode,hidelinks,bookmarks=false]{hyperref}
\newcommand{\ie}{i.e.\ }
\newcommand{\cf}{cf.\ }
\newcommand{\resp}{respectively\ }
\newcommand{\Subsection}{\subsection}
\providecommand{\nobreakdash}{\nobreak\hbox{-}}
\setlength{\emergencystretch}{2em}
\multlinegap0pt
\usepackage{amssymb}
\usepackage{enumerate}
\newtheorem{thm}{Theorem}
\renewcommand\phi{\varphi}
\DeclareMathOperator{\spn}{Span}
\title{Models of holomorphic functions on \newline the symmetrized skew bidisc}
\author{Connor Evans, Zinaida A. Lykova, N. J. Young}
\date{Source: arXiv:2512.02670v2 (CC BY 4.0)}
\begin{document}
\maketitle

\noindent\emph{Source and licence.} Models of holomorphic functions on \newline the symmetrized skew bidisc. Version arXiv:2512.02670v2 (CC BY 4.0), distributed by arXiv under the Creative Commons Attribution 4.0 International licence, http://creativecommons.org/licenses/by/4.0/, as recorded in the arXiv licence metadata for this exact version and shown on the abstract page https://arxiv.org/abs/2512.02670v2. Full licence notice: \texttt{licenses/arxiv-2512.02670-CC-BY-4.0.txt}.\par\smallskip

\noindent\textit{Editorial scope.} Counted unit: the article's thm modelformulaforGr (source0.tex lines 413--428). The article's complete original proof of it is delivered with it (source0.tex lines 429--625, one \texttt{proof} environment of 197 lines). No mathematics is written, repaired or re-derived: the statement and the proof are the article's own lines, in the article's own notation, and no retained line is substituted or re-flowed.  Retained uncounted as the source-local context the proof needs: lines 287--292; lines 372--374; lines 408--411.  Nothing was omitted from any retained span, so the package carries no omission record.  No retained line contains a citation, so the wrapper carries no bibliography and no bibliography entry.  No retained line contains a figure environment, an image-inclusion command or drawing code, so no image is delivered and the package declares no asset dependency.  Editorial formatting only, in addition: an AMS \texttt{amsart} preamble that restates the article's own declarations and macros used by the retained lines, namely the environments \texttt{thm} on separate counters, together with the standard packages those lines need.  The retained environments print in this document as Theorem 1 in order of appearance, whereas the article prints the same items with the numbering of its own full document.  The complete native source of the article as distributed in the arXiv e-print for this exact version is bundled unchanged as \texttt{sources/190211/source0.tex}.
\begin{equation} \label{RonM-intro}
\mathcal{R}=\begin{bmatrix}
	1_{\mathcal{H}_{1}  } 	& 0                \\
	0 		& r\cdot 1_{\mathcal{H}_{1}^{\perp}}  	\\
\end{bmatrix}\in\mathcal{B(M)}.
\end{equation}

\begin{thm}{\normalfont{({\bf{Agler}})}}\label{aglerstheorembidisc}
A function $\phi$ on $\mathbb{D}^{2}$ belongs to the Schur class  $\mathcal{S}(\mathbb{D}^{2})$ if and only if $\phi$ has a model. 
\end{thm}

\begin{equation}\label{Finvolution}
 F(\lambda^{\sigma}) = f(\pi(r\lambda_{2},\lambda_{1})) = f(\lambda_{1}+r\lambda_{2},r\lambda_{1}\lambda_{2}) = F(\lambda), \ \text{for all} \
  \lambda \in  r\mathbb{D}\times\mathbb{D}.
\end{equation}

\begin{thm}\label{modelformulaforGr}
Let $f\in \mathrm{Hol}(\mathbb{G}_{r},\overline{\mathbb{D}})$ and let 
$$F = f \circ \pi \circ T_{r}: \mathbb{D}^{2} \to \overline{\mathbb{D}}.$$ Then 
there exist a complex Hilbert space $\mathcal{M}$,  a closed non-trivial proper subspace  $\mathcal{H}_{1}$ of $\mathcal{M}$,
a unitary operator $U$ on $\mathcal{M}$, 
 a holomorphic map $w: r\mathbb{D}\times \mathbb{D} \to \mathcal{M}$, which satisfies $w(\lambda^{\sigma}) = w(\lambda)$ for all $\lambda \in r\mathbb{D}\times\mathbb{D}$,   such that, for all $\lambda, \mu \in r\mathbb{D}\times\mathbb{D}$, 
\begin{equation}\label{Zrinnerproduct}
1-\overline{F(\mu)}F(\lambda) = \langle Z_{r}(\lambda,\mu)w(\lambda),w(\mu) \rangle_{\mathcal{M}}, 
\end{equation}
where
\begin{align*}
Z_{r}(\lambda,\mu) =& (1_{\mathcal{M}}-r\overline{\mu_{2}}\mathcal{R}^{-1}U^{*})(1_{\mathcal{M}}-  \overline{\mu}_{1}\lambda_{1}\mathcal{R}^{-2})(1_{\mathcal{M}}-r\lambda_{2}U\mathcal{R}^{-1}) \nonumber \\
+ & (1_{\mathcal{M}}-\overline{\mu_{1}}\mathcal{R}^{-1}U^{*})(1_{\mathcal{M}}-r^{2} \overline{\mu}_{2}\lambda_{2}\mathcal{R}^{-2})(1_{\mathcal{M}}-\lambda_{1}U\mathcal{R}^{-1})
\end{align*}
and $\mathcal{R}\in\mathcal{B(M)}$ is defined by equation \eqref{RonM-intro}.
\end{thm}

\begin{proof}
Since $F \in \mathcal{S}(\mathbb{D}^{2})$, by Agler's Theorem \ref{aglerstheorembidisc}, $F$ has a model $(\mathcal{H},u)$, that is, there exists an orthogonally decomposed Hilbert space $\mathcal{H}=\mathcal{H}_{1}\oplus\mathcal{H}_{2}$ and a pair of holomorphic maps $u=(u_{1},u_{2})$ from $\mathbb{D}^{2}$ to $\mathcal{H}_{1},\mathcal{H}_{2}$ respectively such that, for all $\lambda, \mu \in \mathbb{D}^{2}$, 
\begin{equation}\label{bidiscmodel}
1-\overline{F(\mu)}{F(\lambda)}=\langle(1-\overline{\mu_{1}}\lambda_{1}) u_{1}({\lambda}),u_{1}({\mu})\rangle_{\mathcal{H}_{1}}+\langle(1-\overline{\mu}_{2}\lambda_{2}) u_{2}({\lambda}),u_{2}({\mu})\rangle_{\mathcal{H}_{2}}.
\end{equation}

Consider $\lambda$ and $\mu$ in $ r\mathbb{D}\times \mathbb{D}$,
replace $\lambda, \mu$ by $\lambda^{\sigma}, \mu^{\sigma}$ respectively in equation \eqref{bidiscmodel} and use equation \eqref{Finvolution} to deduce that, for all $\lambda$ and $ \mu $ in $ r\mathbb{D}\times \mathbb{D}$, the following equation holds
\begin{align}\label{bidiscmodelsigma}
1-&\overline{F(\mu)}{F(\lambda)} \nonumber\\
&=(1-r^{2}\overline{\mu_{2}}\lambda_{2})\langle{u_{1}({\lambda^{\sigma}}),u_{1}({\mu}^{\sigma})\rangle_{\mathcal{H}_{1}}}+\langle{(1-r^{-2}\overline{\mu_{1}}\lambda_{1})u_{2}({\lambda^{\sigma}}),u_{2}({\mu^{\sigma}})\rangle_{\mathcal{H}_{2}}}.
\end{align}
Take the average of equations (\ref{bidiscmodel}) and (\ref{bidiscmodelsigma}) to obtain, for  all $\lambda$ and $\mu$ in $r\mathbb{D}\times\mathbb{D}$, 
\begin{align*}
1-& \overline{F(\mu)}F(\lambda) \\
= & \dfrac{1}{2}\Bigg(\Big\langle(1-\overline\mu_{1}\lambda_{1})u_{1}({\lambda}),u_{1}(\mu)\Big\rangle_{\mathcal{H}_{1}}+ \Big\langle (1-r^{-2}\overline\mu_{1}\lambda_{1})u_{2}({\lambda^{\sigma}}), u_{2}(\mu^{\sigma})\Big\rangle_{\mathcal{H}_{2}} \\
&+  \Big\langle(1-r^{2}\overline{\mu}_{2}\lambda_{2}) u_{1}(\lambda^{\sigma}), u_{1}(\mu^{\sigma})\Big\rangle_{\mathcal{H}_{1}}+\Big\langle(1-\overline{\mu}_{2}\lambda_{2})u_{2}(\lambda), u_{2}(\mu)\Big\rangle_{\mathcal{H}_{2}}\Bigg). 
\end{align*}
The last equation can be  be re-written as
\begin{align}
1-\overline{F(\mu)}F(\lambda) =& \dfrac{1}{2}\Bigg(\Bigg\langle
\begin{bmatrix}
	(1-\overline\mu_{1}\lambda_{1})u_{1}(\lambda) \\
	(1-r^{-2}\overline\mu_{1}\lambda_{1})u_{2}(\lambda^{\sigma}) \nonumber \\
\end{bmatrix},
\begin{bmatrix}
           u_{1}(\mu) \\
           u_{2}(\mu^{\sigma}) \\
\end{bmatrix}\Bigg\rangle_{\mathcal{H}_{1}\oplus \mathcal{H}_{2}}\\
& + \Bigg\langle
\begin{bmatrix}
	(1-r^{2}\overline\mu_{2}\lambda_{2})u_{1}({\lambda^{\sigma}}) \\
	(1-\overline\mu_{2}\lambda_{2})u_{2}({\lambda}) \\
\end{bmatrix},
\begin{bmatrix}
	u_{1}({\mu}^{\sigma}) \\
	u_{2}({\mu}) \\
\end{bmatrix}\Bigg\rangle_{\mathcal{H}_{1}\oplus \mathcal{H}_{2}}\Bigg). \label{bulkyequation}
\end{align}
For each $\lambda \in r\mathbb{D}\times\mathbb{D}$, define
the vector $v(\lambda)\in\mathcal{H}$ and the operator $\widetilde{\mathcal{R}} \in \mathcal{B(H)}$ by
\begin{equation}\nonumber
v({\lambda}) = \dfrac{1}{\sqrt{2}}
\begin{bmatrix}
	u_{1}({\lambda}) \\
	u_{2}({\lambda^{\sigma}}) \\
\end{bmatrix},~
\widetilde{\mathcal{R}} = 
\begin{bmatrix}
	1_{\mathcal{H}_1}   	& 0                \\
	0 		& r\cdot 1_{\mathcal{H}_2}  	\\
\end{bmatrix}.
\end{equation}
Then, for all $\lambda,\mu \in r\mathbb{D}\times\mathbb{D}$, equation (\ref{bulkyequation}) can be written as
\begin{equation}\label{averagedequations}
1-\overline{F(\mu)}F(\lambda) = \Big\langle(1_{\mathcal{H}}-\overline{\mu}_{1}\lambda_{1}\widetilde{\mathcal{R}}^{-2})v({\lambda}),v({\mu})\Big\rangle_{\mathcal{H}}+\Big\langle(1_{\mathcal{H}}-r^{2}\overline{\mu}_{2}\lambda_{2}\widetilde{\mathcal{R}}^{-2})v({\lambda}^{\sigma}),v(\mu^{\sigma})\Big\rangle_{\mathcal{H}}.
\end{equation}
Again, use the fact that $F(\lambda^{\sigma})=F(\lambda)$ for all $\lambda\in r \mathbb{D}\times\mathbb{D}$ and replace $\lambda$ with $\lambda^{\sigma}$ in equation (\ref{averagedequations}) to obtain
\begin{equation}\label{newaveragedequations}
1-\overline{F(\mu)}F(\lambda) = \Big\langle(1_{\mathcal{H}}-r\overline{\mu}_{1}\lambda_{2}\widetilde{\mathcal{R}}^{-2})v(\lambda^{\sigma}),v(\mu)\Big\rangle_{\mathcal{H}}+\Big\langle(1_{\mathcal{H}}-r\overline{\mu}_{2}\lambda_{1}\widetilde{\mathcal{R}}^{-2})v({\lambda}),v({\mu}^{\sigma})\Big\rangle_{\mathcal{H}}.
\end{equation}
We then equate the right hand sides of equations (\ref{averagedequations}) and (\ref{newaveragedequations}) to see that
\begin{align*}
\Big\langle(1_{\mathcal{H}}-\overline{\mu}_{1}&\lambda_{1}\widetilde{\mathcal{R}}^{-2})v({\lambda}),v({\mu})\Big\rangle_{\mathcal{H}}+\Big\langle(1_{\mathcal{H}}-r^{2}\overline{\mu}_{2}\lambda_{2}\widetilde{\mathcal{R}}^{-2})v({\lambda}^{\sigma}),v({\mu}^{\sigma})\Big\rangle_{\mathcal{H}}\nonumber \\ 
= \Big\langle(1_{\mathcal{H}}-&r\overline{\mu}_{1}\lambda_{2}\widetilde{\mathcal{R}}^{-2})v(\lambda^{\sigma}),v(\mu)\Big\rangle_{\mathcal{H}}+\Big\langle(1_{\mathcal{H}}-r\overline{\mu}_{2}\lambda_{1}\widetilde{\mathcal{R}}^{-2})v({\lambda}),v({\mu}^{\sigma})\Big\rangle_{\mathcal{H}}. 
\end{align*}
Expanding brackets, we find that
\begin{align}
& \Big\langle v(\lambda), v(\mu) \Big\rangle_{\mathcal{H}} -  \Big\langle \overline{\mu}_{1}\lambda_{1} \widetilde{\mathcal{R}}^{-2} v(\lambda), v(\mu)\Big\rangle_{\mathcal{H}} \nonumber\\
& \hspace{1cm} + \Big\langle v(\lambda^{\sigma}), v(\mu^{\sigma})\Big\rangle_{\mathcal{H}} - \Big\langle r^{2}\overline{\mu}_{2}\lambda_{2}\widetilde{\mathcal{R}}^{-2}v(\lambda^{\sigma}), v(\mu^{\sigma}) \Big\rangle_{\mathcal{H}} \nonumber \\
& =\Big\langle v(\lambda^{\sigma}), v(\mu) \Big\rangle_{\mathcal{H}} - \Big\langle r\overline{\mu}_{1}\lambda_{2}\widetilde{\mathcal{R}}^{-2}v(\lambda^{\sigma}), v(\mu) \Big\rangle_{\mathcal{H}}
\nonumber\\
& \hspace{1cm} + \Big\langle v(\lambda), v(\mu^{\sigma})\Big\rangle_{\mathcal{H}}-\Big\langle r\overline{\mu}_{2}\lambda_{1}\widetilde{\mathcal{R}}^{-2}v(\lambda),v(\mu^{\sigma})\Big\rangle_{\mathcal{H}} \label{rearrangingthisequation}.
\end{align}
Rearrange equation (\ref{rearrangingthisequation}) to obtain, for all $\lambda,\mu\in r\mathbb{D}\times \mathbb{D}$,
\begin{align*}
\Big\langle v(\lambda), v(\mu) \Big\rangle_{\mathcal{H}} &+ \Big\langle v(\lambda^{\sigma}),v(\mu^{\sigma})\Big\rangle_{\mathcal{H}}
- \Big\langle v(\lambda^{\sigma}), v(\mu)\Big\rangle_{\mathcal{H}} - \Big\langle v(\lambda), v(\mu^{\sigma})\Big\rangle_{\mathcal{H}} \\
=&\Big\langle \overline{\mu}_{1}\lambda_{1} \widetilde{\mathcal{R}}^{-2} v(\lambda), v(\mu)\Big\rangle_{\mathcal{H}} +\Big\langle r^{2}\overline{\mu}_{2}\lambda_{2}\widetilde{\mathcal{R}}^{-2}v(\lambda^{\sigma}), v(\mu^{\sigma}) \Big\rangle_{\mathcal{H}} \\
&- \Big\langle r\overline{\mu}_{1}\lambda_{2}\widetilde{\mathcal{R}}^{-2}v(\lambda^{\sigma}), v(\mu) \Big\rangle_{\mathcal{H}} -\Big\langle r\overline{\mu}_{2}\lambda_{1}\widetilde{\mathcal{R}}^{-2}v(\lambda),v(\mu^{\sigma})\Big\rangle_{\mathcal{H}}.
\end{align*}
The last  equation can be simplified to 
\begin{align*}
\Big\langle v(\lambda) - v(\lambda^{\sigma}), & \ v(\mu)\Big\rangle_{\mathcal{H}} + \Big\langle v(\lambda^{\sigma})-v(\lambda), v(\mu^{\sigma})\Big\rangle_{\mathcal{H}} \\
 =&\Big\langle \overline{\mu}_{1}\lambda_{1} \widetilde{\mathcal{R}}^{-2} v(\lambda)-r\overline{\mu}_{1}\lambda_{2}\widetilde{\mathcal{R}}^{-2}v(\lambda^{\sigma}), v(\mu)\Big\rangle_{\mathcal{H}} \\
&+ \Big\langle r^{2}\overline{\mu}_{2}\lambda_{2}\widetilde{\mathcal{R}}^{-2}v(\lambda^{\sigma}) - r\overline{\mu}_{2}\lambda_{1}\widetilde{\mathcal{R}}^{-2}v(\lambda),v(\mu^{\sigma}) \Big\rangle_{\mathcal{H}} 
\end{align*}
and then to
\begin{align}\label{right-form-1}
\Big\langle v(\lambda) - v(\lambda^{\sigma}), & \ v(\mu)-v(\mu^{\sigma})\Big\rangle_{\mathcal{H}} \nonumber \\
=&\Big\langle \overline{\mu}_{1}\widetilde{\mathcal{R}}^{-2}\big(\lambda_{1} v(\lambda)-r\lambda_{2}v(\lambda^{\sigma})\big), v(\mu)\Big\rangle_{\mathcal{H}} \nonumber \\
&+\Big\langle r\overline{\mu}_{2}\widetilde{\mathcal{R}}^{-2}\big(r\lambda_{2}v(\lambda^{\sigma}) - \lambda_{1}v(\lambda)\big),v(\mu^{\sigma}) \Big\rangle_{\mathcal{H}} .
\end{align}
The equation \eqref{right-form-1} can then be written in the form
\begin{align*}
\Big\langle v(\lambda) - v(\lambda^{\sigma}),& \ v(\mu)-v(\mu^{\sigma})\Big\rangle_{\mathcal{H}} \\
=&\Big\langle \widetilde{\mathcal{R}}^{-1}\big(\lambda_{1} v(\lambda)-r\lambda_{2}v(\lambda^{\sigma})\big), \mu_{1}\widetilde{\mathcal{R}}^{-1}v(\mu)\Big\rangle_{\mathcal{H}} \\
&+\Big\langle \widetilde{\mathcal{R}}^{-1}\big(r\lambda_{2}v(\lambda^{\sigma}) - \lambda_{1}v(\lambda)\big),r\mu_{2}\widetilde{\mathcal{R}}^{-1}v(\mu^{\sigma}) \Big\rangle_{\mathcal{H}} 
\end{align*}
and further simplified to the form
\begin{align*}
\Big\langle v(\lambda) - v(\lambda^{\sigma}), & \ v(\mu)-v(\mu^{\sigma})\Big\rangle_{\mathcal{H}} \\
=&\Big\langle \widetilde{\mathcal{R}}^{-1}\big(\lambda_{1} v(\lambda)-r\lambda_{2}v(\lambda^{\sigma})\big), \widetilde{\mathcal{R}}^{-1}\big(\mu_{1}v(\mu)-r\mu_{2}v(\mu^{\sigma})\big)\Big\rangle_{\mathcal{H}}.
\end{align*}
This is equivalent to saying that the Gramian of the family 
$$\{v({\lambda})-v({\lambda}^{\sigma}):\lambda\in r\mathbb{D}\times\mathbb{D}\}$$ in $\mathcal{H}$ is equal to the Gramian of the family 
$$\{\widetilde{\mathcal{R}}^{-1}(\lambda_{1}v({\lambda})-r\lambda_{2}v({\lambda}^{\sigma})) : \lambda \in r\mathbb{D}\times\mathbb{D}\},$$ also in $\mathcal{H}$. Hence
there exists a linear isometry 
\begin{align*}
 L :  \overline{\spn}\Big\{\widetilde{\mathcal{R}}^{-1}(\lambda_{1}v({\lambda})- r\lambda_{2}v({\lambda}^{\sigma})) & : \lambda \in r\mathbb{D}\times\mathbb{D}\Big\} \\
& \rightarrow \overline{\spn}\Big\{v({\lambda})-v({\lambda}^{\sigma}) : \lambda\in r\mathbb{D}\times\mathbb{D}\Big\}
\end{align*}
with
\begin{equation}\label{isometryonthegramian}
L\Big( \widetilde{\mathcal{R}}^{-1}(\lambda_{1}v({\lambda})-r\lambda_{2}v({\lambda}^{\sigma})\Big) = v({\lambda})-v({\lambda}^{\sigma}),
\end{equation}
for all $\lambda \in r\mathbb{D}\times \mathbb{D}$. For subsequent calculations, it becomes advantageous to extend $L$ to a unitary operator $U$ on a Hilbert space $\mathcal{M}\supseteq \mathcal{H}$. We also extend $\widetilde{\mathcal{R}}$ to an operator $\mathcal{R}$ on the Hilbert space $\mathcal{M}=\mathcal{H}_{1}\oplus \mathcal{H}_{1}^{\perp}$, where $\mathcal{H}_{1}^{\perp}= \mathcal{M} \ominus \mathcal{H}_{1}$,  by
\begin{equation*}
\mathcal{R} = 
\begin{bmatrix}
	\widetilde{\mathcal{R}}_{\mathcal{H}}   	& 0                \\
	0 		& r_{\mathcal{H}^{\perp}}  	\\
\end{bmatrix} = 
\begin{bmatrix}
	1_{\mathcal{H}_{1}}   	& 0			& 0              \\
	0 		& r\cdot 1_{\mathcal{H}_{2}}		& 0		\\
	0		& 0					& r\cdot 1_{\mathcal{H}^{\perp}}	\\
\end{bmatrix} = 
\begin{bmatrix}
	1_{\mathcal{H}_{1}  } 	& 0                \\
	0 		& r\cdot 1_{\mathcal{H}_{1}^{\perp}}  	\\
\end{bmatrix}.
\end{equation*}
We rearrange equation (\ref{isometryonthegramian}) with $L$ replaced by $U$ to obtain 
\begin{equation}\label{equalitybetweenextension}
(1_{\mathcal{M}}-\lambda_{1}U\mathcal{R}^{-1})v({\lambda}) = (1_{\mathcal{M}}-r\lambda_{2}U\mathcal{R}^{-1})v({\lambda}^{\sigma}),
\end{equation}
for all $\lambda \in r\mathbb{D}\times\mathbb{D}$.
Since $\mathcal{R}$ is a diagonal operator on $\mathcal{M}$ and $\lambda_{1}\in r\mathbb{D}$, we obtain
$$\| \lambda_{1}U\mathcal{R}^{-1}\|_{\mathcal{B(M)}}  = \lvert \lambda_{1} \rvert \| \mathcal{R}^{-1} \|_{\mathcal{B(M)}} = \dfrac{\lvert \lambda_{1} \rvert}{r} < 1. $$
Hence $1_{\mathcal{M}}-\lambda_{1}U\mathcal{R}^{-1}$ is invertible. Likewise, since 
$$ \| r \lambda_{2} U\mathcal{R}^{-1}\|_{\mathcal{B(M)}} =r \lvert \lambda_{2}\rvert \| \mathcal{R}^{-1}\|_{\mathcal{B(M)}} = \lvert \lambda_{2} \rvert < 1,$$ 
the operator $1_{\mathcal{M}}-r\lambda_{2}U\mathcal{R}^{-1}$ is also invertible. Note that 
\begin{equation}\label{verifyinvertibility}
(1_{\mathcal{M}}-r\lambda_{2}U\mathcal{R}^{-1})(1_{\mathcal{M}}-\lambda_{1}U\mathcal{R}^{-1}) = (1_{\mathcal{M}}-\lambda_{1}U\mathcal{R}^{-1})(1_{\mathcal{M}}-r\lambda_{2}U\mathcal{R}^{-1}),
\end{equation}
which can be verified by expanding brackets. Multiply both sides of equation (\ref{verifyinvertibility}) on the left and right by $(1_{\mathcal{M}}-\lambda_{1}U\mathcal{R}^{-1})^{-1}$ to produce
\begin{equation}\label{commutingequations}
(1_{\mathcal{M}}-\lambda_{1}U\mathcal{R}^{-1})^{-1}(1_{\mathcal{M}}-r\lambda_{2}U\mathcal{R}^{-1}) = (1_{\mathcal{M}}-r\lambda_{2}U\mathcal{R}^{-1})(1_{\mathcal{M}}-\lambda_{1}U\mathcal{R}^{-1})^{-1}.
\end{equation}
Rearrange equation (\ref{equalitybetweenextension}) to
\begin{equation}\label{firstinversestep}
v({\lambda}) = (1_{\mathcal{M}}-\lambda_{1}U\mathcal{R}^{-1})^{-1}(1_{\mathcal{M}}-r\lambda_{2}U\mathcal{R}^{-1})v({\lambda}^{\sigma}),
\end{equation}
for all $\lambda \in r\mathbb{D}\times\mathbb{D}$.
By equation (\ref{commutingequations}), equation (\ref{firstinversestep}) can be written
\begin{equation}\label{rearrangednewformula-1}
 v({\lambda}) =  (1_{\mathcal{M}}-r\lambda_{2}U\mathcal{R}^{-1})(1_{\mathcal{M}}-\lambda_{1}U\mathcal{R}^{-1})^{-1}v({\lambda}^{\sigma}).
\end{equation}
Thus
\begin{equation}\label{rearrangednewformula}
 (1_{\mathcal{M}}-r\lambda_{2}U\mathcal{R}^{-1})^{-1}v({\lambda}) = (1_{\mathcal{M}}-\lambda_{1}U\mathcal{R}^{-1})^{-1}v({\lambda}^{\sigma}),
\end{equation}
for all $\lambda \in r\mathbb{D}\times\mathbb{D}$.
Let us define $w: r\mathbb{D}\times\mathbb{D}\rightarrow \mathcal{M}$ by
\begin{equation}\label{wequation}
w(\lambda)=(1_{\mathcal{M}}-r\lambda_{2}U\mathcal{R}^{-1})^{-1}v({\lambda})~~\text{for all}~\lambda \in r\mathbb{D}\times \mathbb{D}.
\end{equation}
Note, for $\lambda \in r \mathbb{D}\times\mathbb{D}$,
\begin{align}
v({\lambda}) &= (1_{\mathcal{M}}-r\lambda_{2}U\mathcal{R}^{-1})w(\lambda) \label{vlambda},\\
~v({\lambda}^{\sigma}) &= (1_{\mathcal{M}}-\lambda_{1}U\mathcal{R}^{-1})w(\lambda^{\sigma}) \label{vlambdasigma}.
\end{align}
Thus, by equation (\ref{rearrangednewformula}), for $\lambda \in r \mathbb{D} \times \mathbb{D}$,
\begin{align}
w({\lambda^{\sigma}}) &= (1_{\mathcal{M}}-\lambda_{1}U\mathcal{R}^{-1})^{-1}v({\lambda}^{\sigma}) \nonumber \\
& =  (1_{\mathcal{M}}-r\lambda_{2}U\mathcal{R}^{-1})^{-1}v({\lambda}) \nonumber \\
& = w(\lambda). \label{wsigma}
\end{align}
Hence $w$ is symmetric with respect to the involution $\sigma$ on $r\mathbb{D}\times\mathbb{D}$. Substituting the expressions (\ref{vlambda}) and (\ref{vlambdasigma}) into equation (\ref{averagedequations}) and enlarging $\mathcal{H}$ to $\mathcal{M}$, we find that, for all $\lambda \in r\mathbb{D}\times\mathbb{D}$,
\begin{align*}
1-&\overline{F(\mu)}F(\lambda) \\
= &\Big\langle(1_{\mathcal{M}}-\overline{\mu}_{1}\lambda_{1}\mathcal{R}^{-2})(1_{\mathcal{M}}-r\lambda_{2}U\mathcal{R}^{-1}) w(\lambda),(1_{\mathcal{M}}-r\mu_{2}U\mathcal{R}^{-1})w(\mu)\Big\rangle_{\mathcal{M}}  \\
&+\Big\langle (1_{\mathcal{M}}-r^{2} \overline{\mu}_{2}\lambda_{2}\mathcal{R}^{-2})(1_{\mathcal{M}}-\lambda_{1}U\mathcal{R}^{-1}) w(\lambda), (1_{\mathcal{M}}-\mu_{1}U\mathcal{R}^{-1}) w(\mu) \Big\rangle_{\mathcal{M}}\\
= & \Big\langle(1_{\mathcal{M}}-r\mu_{2}U\mathcal{R}^{-1})^{*}(1_{\mathcal{M}}-\overline{\mu}_{1}\lambda_{1}\mathcal{R}^{-2})(1_{\mathcal{M}}-r\lambda_{2}U\mathcal{R}^{-1}) w(\lambda),w(\mu)\Big\rangle_{\mathcal{M}}  \\
&+\Big\langle (1_{\mathcal{M}}-\mu_{1}U\mathcal{R}^{-1})^{*}(1_{\mathcal{M}}-r^{2} \overline{\mu}_{2}\lambda_{2}\mathcal{R}^{-2})(1_{\mathcal{M}}-\lambda_{1}U\mathcal{R}^{-1}) w(\lambda), w(\mu) \Big\rangle_{\mathcal{M}}\\
= &\Big\langle(1_{\mathcal{M}}-r\overline{\mu_{2}}\mathcal{R}^{-1}U^{*})(1_{\mathcal{M}}-\overline{\mu}_{1}\lambda_{1}\mathcal{R}^{-2})(1_{\mathcal{M}}-r\lambda_{2}U\mathcal{R}^{-1})w(\lambda),w(\mu)\Big\rangle_{\mathcal{M}}  \\
&+\Big\langle (1_{\mathcal{M}}-\overline{\mu_{1}}\mathcal{R}^{-1}U^{*})(1_{\mathcal{M}}-r^{2} \overline{\mu}_{2}\lambda_{2}\mathcal{R}^{-2})(1_{\mathcal{M}}-\lambda_{1}U\mathcal{R}^{-1}) w(\lambda), w(\mu) \Big\rangle_{\mathcal{M}}.
\end{align*}
Thus, for all $\lambda \in r\mathbb{D}\times\mathbb{D}$,
$$ 1-\overline{F(\mu)}F(\lambda) = \langle Z_{r}(\lambda,\mu) w(\lambda),w(\mu)\rangle_{\mathcal{M}}, $$
where
\begin{align*}
Z_{r}(\lambda,\mu) =& (1_{\mathcal{M}}-r\overline{\mu_{2}}\mathcal{R}^{-1}U^{*})(1_{\mathcal{M}}-\overline{\mu}_{1}\lambda_{1}\mathcal{R}^{-2})(1_{\mathcal{M}}- r\lambda_{2}U\mathcal{R}^{-1}) \\
& + (1_{\mathcal{M}}-\overline{\mu_{1}}\mathcal{R}^{-1}U^{*})(1_{\mathcal{M}}-r^{2} \overline{\mu}_{2}\lambda_{2}\mathcal{R}^{-2})(1_{\mathcal{M}}-\lambda_{1}U\mathcal{R}^{-1}).
\end{align*}
Therefore equation (\ref{Zrinnerproduct}) holds.\end{proof}

\end{document}
